% --------------------------------------------------------------------
%  COLOUR PALETTE
% --------------------------------------------------------------------
%  One palette, defined once under UTP* names, and used by every book,
%  box, diagram and chart. Five accents (blue, teal, green, orange, red)
%  with a pale tint of each for box backgrounds, plus navy for headings,
%  gold for the cover, ink and muted for text, and a sixth accent,
%  purple, for charts and layered diagrams that need one more series.

\definecolor{UTPNavy}{HTML}{18324A}
\definecolor{UTPBlue}{HTML}{246B9E}
\definecolor{UTPTeal}{HTML}{138A80}
\definecolor{UTPGreen}{HTML}{2D7D46}
\definecolor{UTPOrange}{HTML}{C86B24}
\definecolor{UTPRed}{HTML}{B34843}
\definecolor{UTPPurple}{HTML}{6B4FA0}
\definecolor{UTPGold}{HTML}{D4A13E}
\definecolor{UTPInk}{HTML}{263238}
\definecolor{UTPMuted}{HTML}{60717C}
\definecolor{UTPPaleBlue}{HTML}{EEF6FB}
\definecolor{UTPPaleTeal}{HTML}{EEF8F6}
\definecolor{UTPPaleOrange}{HTML}{FFF5EA}
\definecolor{UTPPaleRed}{HTML}{FFF0EF}
\definecolor{UTPPalePurple}{HTML}{F0ECF9}
\definecolor{UTPPaleGray}{HTML}{F4F6F7}
\color{UTPInk}

% --------------------------------------------------------------------
%  PAGE LAYOUT AND TABLES
% --------------------------------------------------------------------
%  Paragraph spacing, list indents, and the table look: alternating
%  row tint, top-aligned paragraph columns, ragged-right cells.
%
%  Two extra column types for tabularx tables:
%      D{width}   a fixed-width, ragged-right, vertically centred column
%      Y          a ragged-right column that fills the remaining width
%
%  Every table in every book has the same shape, the one Algorithms &
%  Data Structures set: a utptable frame around a full-width tabularx,
%  booktabs rules, a tinted header row of \utpth cells (centred, bold,
%  navy), and a first column in bold navy when it names the rows.
%  Cells are vertically centred, so use m{} rather than p{} columns.
%
%      \begin{utptable}
%      \begin{tabularx}{\linewidth}{@{}>{\bfseries\leavevmode\color{UTPNavy}}m{0.22\linewidth}X@{}}
%      \toprule
%      \rowcolor{UTPPaleBlue!92}
%      \utpth{Operation} & \utpth{What it does} \\
%      \midrule
%      Insertion & Add a new element. \\
%      \bottomrule
%      \end{tabularx}
%      \end{utptable}
%
%  A book whose tables carry dense text asks for \utpcompacttables after
%  the input: tabularx tables then set in \footnotesize with tighter
%  padding and vertically centred cells.

\setlength{\parindent}{0pt}
\setlength{\parskip}{0.55em}
\setlist[itemize]{leftmargin=1.45em,itemsep=0.22em,topsep=0.25em}
\setlist[enumerate]{leftmargin=1.75em,itemsep=0.28em,topsep=0.3em}
\setlist[description]{font=\bfseries\color{UTPNavy},itemsep=0.25em,topsep=0.3em}
\renewcommand{\arraystretch}{1.34}
\setlength{\tabcolsep}{6.5pt}
\arrayrulecolor{UTPBlue!45}
% Keep paragraph-style columns ragged and vertically centred. Tables use
% fixed m columns beside flexible X columns; giving X the same centring
% keeps a short cell level with a tall neighbour in a multi-line row.
\renewcommand{\tabularxcolumn}[1]{m{#1}}
\makeatletter
\g@addto@macro\@arrayparboxrestore{\RaggedRight\arraybackslash}
\makeatother
% A header cell: centred, bold and navy, on the tinted header row.
\newcommand{\utpth}[1]{\multicolumn{1}{c}{\bfseries\color{UTPNavy}#1}}
\AtBeginEnvironment{tabular}{\small\setlength{\extrarowheight}{0.6pt}\rowcolors{2}{UTPPaleBlue!72}{white}}
\AtBeginEnvironment{tabularx}{\small\setlength{\extrarowheight}{0.6pt}\rowcolors{2}{UTPPaleBlue!72}{white}}
\AtBeginEnvironment{figure}{\centering}
\newcolumntype{D}[1]{>{\raggedright\arraybackslash}m{#1}}
\newcolumntype{Y}{>{\raggedright\arraybackslash}X}
\newcommand{\utpcompacttables}{%
  \renewcommand{\tabularxcolumn}[1]{m{##1}}%
  \AtBeginEnvironment{tabularx}{%
    \footnotesize
    \setlength{\tabcolsep}{4pt}%
    \renewcommand{\arraystretch}{1.25}%
    \setlength{\extrarowheight}{0.6pt}%
  }%
}
\emergencystretch=1.5em
\hbadness=10000
\hfuzz=8pt

% Page-breaking discipline. Without these a paragraph can leave a single
% line stranded at the foot of one page or the head of the next, which is
% the most visible sign of unattended typesetting.
\widowpenalty=10000          % no first line of a paragraph left alone at a page foot
\clubpenalty=10000           % no last line of a paragraph left alone at a page head
\displaywidowpenalty=10000   % the same around displayed equations
\brokenpenalty=10000         % do not break a page after a hyphenated line
\predisplaypenalty=10000     % keep the line before a display with it
% Pages are built from boxes of many heights, so let the last line float
% up rather than stretching the gaps between boxes to force a flush foot.
\raggedbottom
% Prefer a slightly loose line to an overfull one, but do not hyphenate
% freely -- short sans lines look ragged when over-hyphenated.
\tolerance=1200
\hyphenpenalty=200

% Character protrusion and font expansion. Both are what stop a justified
% sans column from showing visible rivers down the page.
\microtypesetup{protrusion=true,expansion=true,tracking=true}

\pgfplotsset{
  compat=1.18,
  every axis/.append style={
    axis line style={draw=UTPMuted!85,line width=0.55pt},
    tick style={draw=UTPMuted!85},
    tick label style={font=\scriptsize\color{UTPMuted}},
    label style={font=\small\color{UTPNavy}},
    title style={font=\small\bfseries\color{UTPNavy}},
    legend style={font=\scriptsize,draw=UTPBlue!22,fill=white,rounded corners=1pt}
  },
  every axis plot/.append style={line width=1.15pt}
}
\tikzset{every picture/.append style={line cap=round,line join=round}}
